\documentstyle[%


righttag%


,11pt%


,amssymb%


,russian%               remove in Eng.


,izvru%                 Set izven% in Eng.


]{amsart}





%


\newcommand{\rank}{\operatorname{rank}}


\newcommand{\tts}[1]{} 





\theoremstyle{plain}


\theorembodyfont{\it}


\newtheorem{thms}{\hskip\parindent Теорема}[section]





\theoremstyle{definition}


\theorembodyfont{\rm}


\newtheorem{notes}{\hskip\parindent Замечание}[section]


\newtheorem{defns}{\hskip\parindent Определение}[section]





\numberwithin{equation}{section}








%\hoffset=-15mm%                  For Eng.


\setcounter{page}{44}


\god{2002}


\nomer{12~(487)} %                        Remove in Eng.


%\nomer{Vol.\,46, No.\,12}%               Set in Eng.


%\str{\pageref{begin}--\pageref{end}}%  Set in Eng.


\udk{517.852}





\author{\it Р.\,Габасов, Ф.М.\,Кириллова}


% NB:   ^^ remove in Eng.


\title{Стабилизация систем с запаздываниями\\ методами


оптимального управления}


\thanks{Работа выполнена в рамках Государственной


программы фундаментальных исследований ``Математические структуры''


при финансовой поддержке Белорусского республиканского фонда


фундаментальных исследований (проект \No\,Ф02Р-008).}





\date{}


%\pagestyle{myheadings}%         Set in Eng.


\pagestyle{plain} %              Remove in Eng.


\begin{document}


%\thispagestyle{plain}%          Set in Eng.


\maketitle


\label{begin}











Статья посвящена изложению одного подхода к решению базовой задачи


теории управления системами с запаздыванием


--- синтезу стабилизирующих обратных связей. Подход


опирается на конструктивную теорию оптимального управления [1],


разработанную авторами и их коллегами. Стабилизация


динамических систем была первой задачей, на основе которой стала


строиться теория управления (регулирования). До начала 50-х гг. XX


века она оставалась центральной задачей теории управления. С


появлением теории оптимального управления интерес к проблеме


стабилизации стал менее заметным, но в последние годы она вновь


привлекает внимание многих ученых, поскольку практическое значение


проблемы стабилизации не уменьшилось и появились новые возможности


решения классической задачи с помощью методов, созданных


в теории управления за последние пятьдесят лет. Системы с


запаздыванием представляют непосредственное обобщение


обыкновенных динамических систем. Современный этап их развития


начался в конце 40-х гг. ХХ века и почти совпал с началом


современного этапа теории управления. Несмотря на значительные


усилия ученых, многие задачи управления остаются для систем с


запаздыванием до сих пор мало исследованными. К таким задачам


относится и задача стабилизации систем с запаздыванием с помощью


ограниченных обратных связей [2], [3].


В данной статье сначала излагается метод решения задачи


стабилизации динамических систем с запаздыванием в состоянии и


управлении. Метод основан на позиционных решениях вспомогательных 


задач оптимального управления. Затем исследуется


проблема синтеза наблюдателей систем с запаздыванием


и после этого решается задача стабилизации систем с запаздыванием


в классе обратных связей по выходу.





\section{Ограниченные стабилизирующие обратные связи по состоянию}





{\bf 1.1.} Рассмотрим динамическую систему, поведение которой при


$t \ge t_*$ описывается уравнением


\begin{equation}\label{11}


x^{(n)}(t)+\sum_{i=0}^{n-1}a_i(t)x^{i}(t)+


\sum_{i=0}^{n-1}a_{1i}(t)x^{(i)}(t-\beta)=u(t).


\end{equation}


Здесь $x(t)$ --- скаляр, характеризующий положение системы в


момент времени $t$; $x^{(i)}(t)$ --- $i$-я производная; $u(t)$,


$t\ge t_*$, --- скалярная кусочно-непрерывная функция (управление);


$a_i(t)$, $a_{1i}(t)$, $i=\overline{0,n-1}$, $t \ge t_{*}$,


--- ограниченные кусочно-непрерывные скалярные функции; $\beta$, $0 < \beta


< \infty$, --- запаздывание.





Как известно, в теории систем с запаздыванием (\ref{11}) отрезок


непрерывной $n$-мерной кривой


$$


x_t(\cdot)=(x^{i}(s)),\ \ s \in


[t-\beta,t],\ \ i=\overline{0,n-1},


$$


называется состоянием системы


(\ref{11}) в момент $t$. Пусть $X_{t}(\cdot)$ --- множество состояний


в момент~$t$.





\begin{defns}


Функционал


\begin{equation}\label{12}


u=u(t,x_t(\cdot)),\ \ x_{t}(\cdot) \in X_t(\cdot),


\ \ t \ge t_*,\ \ u(t,0)=0,\ \ t \ge t_*,


\end{equation}


называется ограниченной стабилизирующей обратной связью (по


состоянию), если при заданном $L$, $0 < L < \infty$,





1) справедливо неравенство


\begin{equation}\label{13}


|u(t,x_t(\cdot))| \le L, x_t(\cdot) \in X_{t}(\cdot),\ \ t \ge t_*;


\end{equation}





2) нулевое решение $x(t)=0$, $t \ge t_*$, замкнутой системы


\begin{equation}\label{14}


x^{(n)}(t)+\sum_{i=0}^{n-1}a_i(t)x^{(i)}(t)+


\sum_{i=0}^{n-1}a_{1i}(t)x^{(i)}(t-\beta)=u(t,x_t(\cdot))


\end{equation}


асимптотически устойчиво в $X_{t_*}(\cdot)$.


\end{defns}





Синтез эффективных (с большой областью притяжения $X_{t_*}(\cdot)$


и определенным качеством переходных процессов в (\ref{14}))


ограниченных стабилизирующих обратных связей (\ref{12}),


(\ref{13}) представляет весьма серьезную математическую проблему.


Ее решение вряд ли возможно без использования современных


вычислительных устройств дискретного действия. В связи с этим в


данной статье задача стабилизации системы (\ref{11}) решается в


классе дискретных обратных связей.





Пусть $h=\beta/M>0$ --- период квантования времени (вообще говоря,


малое число), $T_h=\{t_*, t_*+h, t_*+2h, \dotsc\}.$





\begin{defns}


Функционал


\begin{equation}\label{15}


u=u(t,x_t(\cdot)),\ \ x_t(\cdot) \in X_{t}(\cdot),\ \ t \in T_h,


\end{equation}


называется дискретной обратной связью (с периодом квантования


$h$), если траектория замкнутой нелинейной системы (\ref{14})


совпадает с траекторией линейной системы (\ref{11}) с тем же


начальным состоянием и управлением


$$


u(t)=u(t_*+kh, x_{t_*+kh}(\cdot)),\ \ t \in [t_*+kh, t_*+(k+1)h], \tts{[}


\ \ k=0,1,2,\dotsc


$$


\end{defns}





Проблема аналитического конструирования


ограниченных дискретных стабилизирующих


обратных связей (\ref{15}) не проще аналогичной


проблемы, поставленной для (\ref{12}). Идея описываемой


ниже реализации обратной связи (\ref{15}) в режиме реального


времени идентична той, что предложена в [4]--[6] для обыкновенных


систем. В каждый текущий момент $\tau \in T_h$, $\tau > t_*$, по


доступному состоянию $x_t(\cdot)$ строится оптимальное программное


управление $u^0(t|\tau,x_{\tau}(\cdot))$, $t \in T(\tau)$,


вспомогательной (сопровождающей) задачи оптимального управления


(см.~п.\,1.2). При этом сопровождающая задача не решается заново, а


двойственным методом корректируется решение, полученное в


предыдущий момент $(\tau-h)$. Как показано в [7], эта процедура


чрезвычайно эффективна --- время на ее выполнение составляет


несколько процентов от времени, за которое используемый


микропроцессор интегрирует систему (\ref{11}) на промежутке


$T(\tau).$ Если время коррекции не превосходит $h$ единиц


реального времени, в котором записано уравнение (\ref{11}), то


можно говорить о реализации обратной связи (\ref{15}) в режиме


реального времени. На промежутке времени $[\tau,\tau+h]$ \tts{[} в систему


(\ref{14}) подается сигнал


$u^{*}(t)=u^0(\tau\mid\tau,x^*_\tau,x^*_\tau(\cdot))=


u^0(\tau,x^*_\tau(\cdot))=u(\tau,x_\tau(\cdot))$,


$t \in [\tau, \tau+h]$. \tts{[}


В результате система (\ref{14}) в момент $\tau+h$ оказывается в


состоянии $x^*_{\tau+h}(\cdot),$ и описанная процедура реализации


обратной связи (\ref{15}) повторяется. Для начального момента


$\tau=t_*$ разработана специальная процедура [7].





С точки зрения приложений вполне достаточно вместо


асимптотического свойства $x(t) \rightarrow 0$, $t \rightarrow


\infty$, обеспечить свойство $\|x(t)\| \rightarrow o(h)$, $t


\rightarrow \infty$, естественное при дискретизации управлений.


\medskip





{\bf 1.2.} Пусть $\theta$, $0 < \theta=Nh < \infty$, --- некоторое


число (параметр метода), $u(t)$, $t \ge t_*$, --- дискретное


программное управление: $u(t)=u(t_*+kh)$, $t \in [t_*+kh, (k+1)h]$, \tts{[}


$k=0,1,2,\dotsc$; $x_{\tau}^{*}(\cdot)$ --- состояние системы


(\ref{14}) в текущий момент $\tau \in T_h$. В классе дискретных


управлений $u(t)$, $t \in T(\tau)$, и кусочно-непрерывных управлений


$u(t)$, $t \in [\tau+\theta-\beta, \tau+\theta]$, \tts{[} рассмотрим задачу


оптимального управления, сопровождающую задачу стабилизации


системы (\ref{11}),


\begin{gather}


B(\tau,x^*_{\tau}(\cdot)) =\min


\int_{\tau}^{\tau+\theta+\beta}|u(t)|dt, \notag \\


%


\label{16}


x^{(n)}(t)+\sum_{i=0}^{n-1}a_i(t)x^{(i)}(t)+


\sum_{i=0}^{n-1}a_{1i}(t)x^{(i)}(t-\beta)=u(t), \\


%


x_{\tau}(\cdot)=x_{\tau}^{*}(\cdot),\quad


u(t)=\sum_{i=0}^{n-1}a_{1i}(t)x^{(i)}(t-\beta),\quad t \in


[\tau+\theta, \tau+\theta+\beta],\notag\\


%


x^{(i)}(\tau+\theta)=0,\ \ i =\overline{0,n-1};\ \ |u(t)| \le L,


\ \ t \in T(\tau)=[\tau, \tau+\theta+\beta].\notag


\end{gather}





Быстрые алгоритмы решения задач типа (\ref{16}) для обыкновенных


систем предложены в [7]. Они допускают обобщение и на системы с


запаздыванием. Обозначим через $u^0(t\mid\tau, x_{\tau}^*(\cdot))$, $t


\in T(\tau)$, --- оптимальное программное


управление задачи (\ref{16}), $X_{\tau}^{0}$ ---


множество состояний $x_{\tau}^{*}(\cdot)$, для которых задача


(\ref{16}) имеет решение.





\begin{defns}


Функционал


\begin{equation}\label{17}


u^{0}(\tau, x_{\tau}(\cdot)) = u^{0}(\tau\mid\tau, x_{\tau}(\cdot)),


\ \ x_{\tau}(\cdot) \in X_{\tau}^{\theta}(\cdot),\ \ \tau \in T_{h},


\end{equation}


называется оптимальным стартовым управлением типа дискретной


обратной связи задачи (\ref{16}).


\end{defns}





\begin{thms}


При достаточно малых $h>0$ обратная связь


$$


u(t, x_{t}(\cdot))=u^0(t,x_{t}(\cdot)),


\ \ x_{t}(\cdot) \in X_{t}^{\theta}(\cdot),\ \ t \in T_{h},


$$


является ограниченной стабилизирующей дискретной обратной связью


для системы $(\ref{11})$.


\end{thms}





Доказательство теоремы в основных чертах повторяет доказательство


аналогичного утверждения из [6]. Пусть замкнутая система в момент


$\tau$ находится в состоянии $x_{\tau}(\cdot)$. Значение функции


Беллмана задачи (\ref{16}) на позиции $(\tau, x_{\tau}(\cdot))$


равно $B(\tau, x_{\tau}(\cdot))$. На промежутке $[\tau, \tau+h]$ \tts{[}


подадим в систему (\ref{16}) управление


$u^*(t)=u^0(\tau,x_{\tau}(\cdot))$, $t \in [\tau, \tau+h]$, \tts{[} после


чего она в момент $\tau+h$ окажется в состоянии


$x^*_{\tau+h}(\cdot)$. Из ограничений задачи (\ref{16}) видно, что


управление $u(t)=u^0(t\mid\tau,x_{\tau}^*(\cdot))$, $t \in T(\tau)$;


$u(t)=u^0(t-\beta|\tau,x^*(\tau))$, $t \in [\tau+\theta,


\tau+\theta+\beta]$, \tts{[} порождает траекторию $x^0(t)$, $t \in [\tau,


\tau+\theta+\beta]$ со свойством $x^0(t) \equiv 0$, $t \in


[\tau+\theta, \tau+\theta+\beta]$. Для позиции $(\tau+h,


x_{\tau+h}^*(\cdot))$ рассмотрим управление


$u(t)=u^0(t\mid\tau,x^*_{\tau}(\cdot))$,


$t \in [\tau+h, \tau+\theta]$; \tts{[}


$u(\tau+\theta)=u_h^0(\tau+\theta-\beta)$; $u(t)=u^0(t+\beta)$, $t \in


[\tau+\theta+h, \tau+\theta+\beta]$; \tts{[} $u(\tau+\theta+\beta)=0$, где


$u^0_h(\tau+\theta-\beta)$ --- такое число, что


$|u^0_h(\tau+\theta-\beta)|\cdot h =


\displaystyle\int_{\tau+\theta-\beta}^{\tau+\theta-\beta+h}|u^0(t)|dt$. Это


управление почти допустимо, на нем выполняются все ограничения


задачи (\ref{16}), кроме $x(\tau+\theta+h)=0$. Последнее


выполняется с точностью $o(h)$: $\|x(\tau+\theta+h)\|=o(h)$. На


рассматриваемом управлении критерий качества задачи (\ref{16})


равен $B(\tau,x_{\tau}^*(\cdot))$. Поэтому в худшем случае будем


иметь $B(\tau+h, x_{\tau+h}^*(\cdot)) \le B(\tau,


x^*(\tau))+o(h)$. При достаточно малых $h>0$ найдется такое число


$q$ (не зависящее от $\tau$), при котором $B(\tau+qh, x^*_{\tau+qh


}(\cdot)) < B(\tau,x_{\tau}(\cdot))$. Таким образом, функция


Беллмана строго убывает в моменты $t_*+qh$, $q=1,2,\dotsc$, если


$\|x_{\tau}(\cdot)\| \ne o(h)$. Поскольку она при весьма общих


предположениях определенно положительна и обладает бесконечно


малым высшим пределом, то $x(t_{*}+qh) \rightarrow o(h)$, $q =


1,2,\dotsc$ Используя ограниченность управления, нетрудно


показать, что $\|x(t)\| \rightarrow o(h)$, $t \rightarrow \infty$.





Дополнительные свойства обратной связи (\ref{17}), характеризующие


размеры области притяжения $X_{t_*}^{\theta}(\cdot)$ и качество


переходных процессов в (\ref{14}), аналогичны тем, которые


приведены в [5], [6] для обыкновенных систем.





Реализация обратной связи (\ref{17}) с помощью регуляторов


аналогична описанной в [6].


\medskip





{\bf 1.3.} Рассмотрим теперь более общую систему с запаздыванием


\begin{equation}\label{18}


\Dot x(t)=A(t)x(t)+A_1(t)x(t-\beta)+B(t)u(t),


\end{equation}


в которой $x(t)$, $u(t)$ --- $n$-векторы, $A(t)$, $A_{1}(t)$, $B(t)$,


$t\ge t_*$, --- ограниченные кусочно-непрерывные $n \times n$-матричные


функции, $|\det \ B(t)| > \alpha > 0$, $t \ge t_*$.





Состояние системы (\ref{18}) имеет вид $x_t(\cdot)=x(s)$,


$t-\beta$ $\le s \le t$. Понятия стабилизирующей


ограниченной дискретной обратной связи


вводятся аналогично предыдущему случаю.





Сопровождающая задача оптимального управления теперь примет вид


\begin{gather*}


B(\tau,x_{\tau}^*(\cdot))=\min_{u}\int_{\tau}^{\tau+\theta+\beta}


\|u(t)\|dt, \\


%


\Dot x(t)=A(t)x(t)+A_1(t)x(t-\beta)+B(t)u(t),


\ \ x_{\tau}(\cdot)=x_{\tau}^*(\cdot), \\


%


u(t)=B^{-1}(t)A(t)x(t-\beta),\ \ t \in [\tau+\theta,


\tau+\theta+\beta];\ \ x(\tau+\theta)=0, \\


%


\|u(t)\| \le L,\ \ t \in T(\tau)=[\tau, \tau+\theta+\beta].


\end{gather*}


Здесь $x_{\tau}^{*}(\cdot)$ --- состояние системы (\ref{18}) в


момент $\tau$; $u(t)$, $t \in [\tau, \tau+\theta]$, --- дискретное


управление; $u(t)$, $t \in [\tau+\theta, \tau+\theta+\beta]$, ---


кусочно-непрерывное управление.





Исследование задачи построения ограниченной стабилизирующей


дискретной обратной связи проводится как в п.~1.2.


\medskip





{\bf 1.4.} Перейдем к системе с запаздыванием, стабилизация


которой сложнее предыдущих задач. Рассмотрим систему


\begin{equation}\label{19}


\Dot x (t)=A(t)x(t)+A_1(t)x(t-\beta)+b(t)u(t).


\end{equation}


Здесь $u(t)$ --- скаляр, значение управления в момент $t$, $b(t)$,


$t\ge t_*$, --- ограниченная кусочно-непрерывная $n$-вектор-функция,


остальные элементы обладают теми же свойствами, что и в (1.8).





Сложность задачи стабилизации системы (\ref{19}) определяется


сложностью процедуры управления такими системами, в которых на $n$


переменных $x_{1}(t), \dotsc, x_{n}(t)$ приходится только одна


переменная управления $u(t)$.





Сопровождающая задача теперь имеет вид


\begin{gather}


B(\tau,x_{\tau}^*(\cdot))=\min_{u}\int_{\tau}^{\tau+\theta+\beta}


|u(t)|dt, \label{110}\\


%


\Dot x(t)=A(t)x(t)+A_1(t)x(t-\beta)+b(t)u(t),


\ \ x_{\tau}(\cdot)=x_{\tau}^*(\cdot),\notag \\


%


A_1(t)x(t-\beta)+b(t)u(t)=0,


\ \ t \in [\tau+\theta,\tau+\theta+\beta]; \notag \\


%


x(\tau+\theta)=0,\ \


|u(t)| \le L,\ \ t \in T(\tau)=[\tau, \tau+\theta+\beta].\notag


\end{gather}


Здесь $u(t)$, $t \in [\tau, \tau+\theta]$, \tts{[} --- дискретное управление;


$u(t)$, $t \in [\tau+\theta, \tau+\theta+\beta]$, ---


кусочно-непрерывное управление.





\begin{notes}


Сложность задачи (\ref{110}) определяется


ограничением на $[\tau+\theta,\tau+\theta+\beta]$.


При фактическом решении разумно заменить его


на $-e\varepsilon \le A_1(t)x(t-\beta)+b(t)u(t) \le e\varepsilon$,


$t \in [\tau+\theta, \tau+\theta\beta]$, где $e=(1, \dotsc, 1) \in


R^n$, $\varepsilon > 0$ --- достаточно малое положительное число.


\end{notes}





{\bf 1.5.} Рассмотрим, наконец, динамическую систему с


запаздыванием в управлении


$$


%\begin{equation}\label{111}


\Dot x (t) =A(t)x(t) +b(t)u(t-\beta),\ \ x(t) \in R^n,\ \ u(t) \in R.


%\end{equation}


$$





Пусть $u_t(\cdot)=(u(s)$, $t -\beta \le s<t)$,


$U_t(\cdot)=\{u_t(\cdot):|u(s)| \le L$, $t -\beta \le s \le t\}$.


Состоянием системы (\ref{110}) в момент $t$ будем называть пару


$(x(t), u_t(\cdot))$. Тогда $(t, x(t), u_t(\cdot))$ --- позиция в


момент $t$; $x(t) \in X_t \subset R^n$.





\begin{defns}


Функционал


$$


u=u(t, x(t), u(\cdot)),\ \ x(t) \in X_t,\ \ u_t(\cdot) \in U_t(\cdot),


\ \ t\ge t_*,\ \ u(t,0)=0,\ \ t \ge t_*,


$$


называется стабилизирующей обратной связью, если нулевое решение


уравнения


\begin{equation}\label{112}


\Dot x(t)=A(t)x(t)+b(t)u(t, x(t), u_{t}(\cdot))


\end{equation}


асимптотически устойчиво в $X_{t_*}$ для всех $u_{t_*}(\cdot)


\in U_{t_*}(\cdot)$.


\end{defns}





Понятия ограниченной дискретной обратной связи вводятся по


аналогии с предыдущими случаями.





С целью реализации ограниченных стабилизирующих дискретных


обратных связей рассмотрим в классе дискретных управлений $u(t)$,


$t\in T(\tau)$, следующую сопровождающую задачу оптимального


управления:


\begin{gather}


B(\tau, x^*(\tau), u_{\tau}(\cdot))=\min


\int_{\tau}^{\tau+\theta-\beta}|u(t)|dt, \notag \\


%


\label{113}


\Dot x(t)=A(t)x(t)+b(t)u(t-\beta),\ \ x(\tau)=x^*(\tau),


\ \ u_{\tau}(\cdot)=u^*_{\tau}(\cdot), \\


%


x(\tau+\theta)=0,\ \ u(t)=0,\ \ t \in [t+\theta-\beta, \tau+\theta],\notag\\


%


|u(t)| \le L,\ \ t \in T(\tau)=[\tau,\tau+\theta].\notag


\end{gather}





Стабилизирующая обратная связь, построенная по оптимальному


стартовому управлению типа дискретной связи задачи (\ref{113}),


обеспечивает асимптотическую устойчивость нулевого решения


замкнутой системы (\ref{112}) при $x(t) \rightarrow 0$, $t \rightarrow


\infty$.


\medskip





{\bf 1.6.} Выше рассматривался только один тип сопровождающих


задач оптимального управления. При использовании других типов


задач оптимального управления можно строить стабилизирующие


обратные связи, обеспечивающие другие качественные свойства


переходных процессов в замкнутых системах.








\section{Наблюдатели систем с запаздыванием}





В современной теории управления наблюдатели динамических систем


используются при синтезе обратных связей по выходу. Они


преобразуют доступные сигналы измерительных устройств в оценки


состояний динамических систем, на базе которых регуляторы


вырабатывают управляющие воздействия. С неуклонным ростом


размерности математических моделей исследуемых в наши дни


физических систем значение наблюдателей возрастает, поскольку, как


правило, с повышением точности моделирования количество доступных


выходных сигналов не меняется. В [8] рассмотрена задача синтеза


нелинейных наблюдателей обыкновенных динамических систем. Ниже


аналогичная задача решается для систем с запаздыванием.


\medskip





{\bf 2.1.} Пусть на промежутке времени $t \ge t_{*}$ поведение


динамической системы описывается уравнением


\begin{equation}\label{21}


\Dot x (t)=A(t)x(t)+A_1(t)x(t-\beta),


\end{equation}


где, как и выше, $x(t)$ --- $n$-вектор, $A(t)$, $A_1(t)$, $t \ge


t_*$, --- ограниченные кусочно-непрерывные $n \times n$-матричные


функции, $\beta$, $0 < \beta < \infty$, --- запаздывание.





Будем считать, что вся доступная информация о поведении системы


(\ref{21}) исчерпывается заданием ограниченного множества


$X_{t_*}(\cdot)$ возможных начальных состояний


$x_{t_*}(\cdot)=x(s)$, $t_*-\beta \le s \le t_*$, и сигналом $y(t)$,


$t \ge t_*$, измерительного устройства


\begin{equation}\label{22}


y(t)=c'(t)x(t)


\end{equation}


с непрерывной $n$-вектор-функцией $c(t)$, $t \ge t_*$.





Пусть $\tau$, $\tau \ge t_*$, --- текущий момент времени, 


$x_{\tau}(\cdot)$ --- текущее состояние системы (\ref{21}),


$y_{\tau}[\cdot]=y(s)$, $t_* \le s \le \tau$, --- сигнал устройства


(\ref{22}), записанный к моменту $\tau$, $Y_{\tau}[\cdot]$ ---


множество возможных сигналов $y_{\tau}[\cdot]$.





\begin{defns}


Динамическую систему


\begin{equation}\label{23}


\Dot z (t)=f(t,z_t(\cdot), y_t[\cdot]),\ \ z_{t_*}=\varphi (\cdot),


\end{equation}


назовем наблюдателем системы (\ref{21}), если для любых


$y_{t}[\cdot] \in Y_{t}[\cdot]$, $t \ge t_*$, решения $x_t(\cdot)$,


$z_t(\cdot)$, $t \ge t_*$, уравнений (\ref{21}), (\ref{23})


удовлетворяют соотношению


\begin{equation}\label{24}


\|z_t(\cdot)-x_{t}(\cdot)\| \rightarrow 0,\ \ t \rightarrow \infty,


\ \ \|x_{t}(\cdot)\|=\max_{t-\beta \le s \le t}\|x(s)\|.


\end{equation}


\end{defns}





Наблюдатели (\ref{23}) будем строить с помощью системы управления


\begin{equation}\label{25}


\Dot z(t)=A(t)z(t)+A_1(t)z(t-\beta)+B(t)u(t),


\ \ z_{t_*}(\cdot)=\varphi_{t_*}(\cdot),


\end{equation}


в которой $B(t)$, $t \ge t_{*}$, --- ограниченная кусочно-непрерывная


$n \times n$-матричная функция, $u(t)$ --- значение $n$-вектора 


управления в момент $t$; $\det B(t) \ge \alpha >0$, $t \ge t_*$.





В терминах системы управления (\ref{25}) задача построения


наблюдателя (\ref{23}) сводится к синтезу обратной связи


\begin{equation}\label{26}


u=u(t,z_t(\cdot), y_t[\cdot]),\ \ z_{t}(\cdot) \in Z_{t}(\cdot),


\ \ y_{t}[\cdot] \in Y_{t}[\cdot], t \ge t_{*},


\end{equation}


и начального состояния $\varphi_{t_*}(\cdot)$, при которых


замкнутая система


\begin{equation}\label{27}


\Dot z(t) =A(t)z(t)+A_1(t)z(t-\beta)+B(t)u(t,z_{t}(\cdot),


\, y_t[\cdot]),\ \  z_{t_*}(\cdot)=\varphi_{t_*},


\end{equation}


является наблюдателем системы (\ref{21}).





Обратная связь (\ref{26}) определена на оси непрерывного времени


$t \ge t_*$, т.\,е. может менять свои значения в любой момент


времени $t \ge t_*$. Как и выше, будем строить дискретные обратные


связи, которые могут менять свои значения только в дискретные


моменты времени $t \in T_{h}=\{t_*,t^*+h,t^*+2h, \dotsc\}$, где


$h>0$ --- период квантования времени. При достаточно малых $h$ (что


предполагается в дальнейшем) качество дискретных обратных связей


мало отличается от качества непрерывных аналогов.





При синтезе обратных связей большое значение имеет геометрическое


ограничение


\begin{equation}\label{28}


\|u(t, x_t(\cdot), y_t[\cdot])\| \le L,\ \ z_t(\cdot) \in


Z_{t}(\cdot),\ \ y _{t}[\cdot] \in Y_t[\cdot], t \in T_h.


\end{equation}





Ниже обратные дискретные связи будем строить с учетом ограничения


(\ref{28}).





Траектория нелинейной замкнутой системы (\ref{27}), полученной


после замыкания системы (\ref{21}) дискретной обратной связью


(\ref{26}), представляет траекторию линейной системы (\ref{21}) с


начальным условием $z_{t_*}(\cdot) = \varphi_{t_*}(\cdot)$ и


управлением


$$


u(t)=u(t_*+kh, z_{t_*+kh}(\cdot), y_{t_{*}+kh}[\cdot]),\ \ t \in


[t_*+kh, t_*+(k+1)h],\tts{[} \ \ k=0,1,2, \dotsc


$$





После перехода к дискретным обратным связям естественно


соотношение (\ref{24}) заменить на следующее:


$$


\|z_t(\cdot) - x_t(\cdot)\| \rightarrow o(h),\ \ t \rightarrow\infty.


$$





Пусть $\tau \in T_h$ --- текущий момент времени, $\tau \ge (M+1)h$,


$Mh=\beta$; $z_{\tau}^*(\cdot)$ --- текущее состояние наблюдателя


(\ref{27}), $y_{\tau}[\cdot]$ --- сигнал устройства (\ref{22}),


записанный к моменту $\tau$, $0=\mu_0(\tau)<


\mu_1(\tau)<\dotsb<\mu_{Mn}(\tau)$ --- некоторые числа, кратные $h$;


выполняется условие наблюдаемости


$$


\rank \left(


\begin{array}{c}


c'(\tau-\mu_{i}(\tau))F(\tau-\mu_i(\tau), t_*),


\ \ \displaystyle\int_{t_*-(k+1)h}^{t_*-kh}c'(\tau-\mu_i(\tau))


\times \\[3mm] \times F(\tau-\mu_i(\tau),s+\beta)


A_1(s+\beta)ds,\ \ k=\overline{0,M-1},\ \ i=\overline{0,Mn}


\end{array} \right) =(M+1)n,


$$


где $F(t,s)$, $t \ge s$, --- фундаментальная $n \times n$-матрица


решений уравнения (\ref{21}); $x_{k}(\tau) \in R^n$,


$k=\overline{-1,M-1}$, --- решение системы уравнений


\begin{gather*}


c'(\tau-\mu_i(\tau))F(\tau-\mu_i(\tau),


t_*)x_{-1}+\sum_{k=0}^{M-1}\int_{t_*-(k+1)h}^{t_*-kh}


c'(\tau-\mu_i(\tau))\times \\


%


\times F(\tau-\mu_i(\tau),s+\beta)A_{1}(s+\beta)ds \;


x_k=y(\tau-\mu_{i}(\tau)),\ \ i=\overline{0,Mn}; \\


%


\wh x(t)=F(t,t_*)x_{-1}(\tau)+\sum^{M-1}_{k=0}


\int_{t_*-(k+1)h}^{t_*-kh}


F(t,s+\beta)A_1(s+\beta)ds \; x_{k}(\tau),\ \ t \ge t_*.


\end{gather*}





В классе дискретных управлений $u(t)$, $t \in [\tau, \tau+\theta]$, \tts{[}


и кусочно-непрерывных управлений $u(t)$, $t \in [\tau+\theta,


\tau+\theta+\beta]$, рассмотрим вспомогательную (сопровождающую)


задачу оптимального управления


\begin{gather}


B(\tau,z_{\tau}^*(\cdot), y_\tau[\cdot]) = \min


\int_{\tau}^{\tau+\theta+\beta} \|u(t)\| dt, \notag \\


%


\label{29}


\Dot z(t)=A(t)z(t)+A_1(t)z(t-\beta)+B(t)u(t),


\ \ z_{\tau}(\cdot)=z_{\tau}^*(\cdot), \\


%


B(t)u(t)+A_{1}(t)z(t-\beta)-A_1(t)\wh x(t-\beta)=0,\ \ t \in


[\tau+\theta, \tau+\theta+\beta], \notag \\


%


z(\tau+\theta)=\wh x(\tau+\theta),\ \ \|u(t)\| \le L,\ \ t \in


[\tau,\tau+\theta+\beta].\notag


\end{gather}


Здесь $B(t)$, $t \in [\tau,\tau+\theta+\beta]$, --- ограниченная


кусочно-непрерывная $n \times n$-матричная функция,


обеспечивающая управляемость системы (\ref{29}), $\det B(t)\ge


\alpha > 0$, $t \ge t_*$, $u(t)$ --- $n$-вектор управления,


$\theta=Nh$ --- параметр метода. Обзначим через


$u^0(t|\tau,z_\tau^*(\cdot), y_{\tau}[\cdot])$, $t \in


[\tau,\tau+\theta+\beta]$, оптимальное программное управление


задачи (\ref{29}), $Z_{\tau}(\cdot)$ --- множество всех состояний,


для которых задача (\ref{29}) имеет решение для всех


$y_{\tau}[\cdot] \in Y_{\tau}[\cdot]$.





\begin{defns}


Функция


\begin{gather*}


u^0(\tau,z_{\tau}(\cdot),


y_{\tau}[\cdot])=u^0(\tau\mid\tau,z_{\tau}(\cdot),y_{\tau}[\cdot]), \\


%


z_{\tau}(\cdot) \in Z_{\tau}(\cdot),\ \ y_{\tau}[\cdot] \in


Y_{\tau}[\cdot],\ \ \tau \in T_{h},


\end{gather*}


называется оптимальным стартовым управлением типа дискретной


обратной связи задачи (\ref{29}).


\end{defns}





Для решения задачи (\ref{29}) можно обобщить быстрые алгоритмы,


разработанные в [7], [9] для обыкновенных систем.





\begin{thms}


Обратная связь


\begin{gather*}


u(t,z_{t}(\cdot), y_{t}[\cdot])=u^0(t,z_{t}(\cdot),y_{t}[\cdot]), \\


%


z_{t}(\cdot) \in Z_{t}(\cdot),\ \ y_{t}[\cdot] \in Y_{t}[\cdot],


\ \ t\in T_{h},


\end{gather*}


и состояние $z_{t_*}(\cdot) \in Z_{t_*}(\cdot)$ синтезируют


наблюдатель $(\ref{27})$.


\end{thms}





Доказательство теоремы 2.1, с одной стороны, основано на


доказательстве аналогичной теоремы для обыкновенных наблюдателей


[8], а с другой --- на доказательстве стабилизирующих свойств


оптимальных обратных связей для систем с запаздыванием (см.


п.~1.1).





\begin{notes}


Начальное состояние $\varphi_{t_*}(\cdot)$


системы (\ref{21}) считалось произвольной непрерывной


(кусочно-непрерывной) функцией из $X_{t_*}(\cdot)$. Часто оно


принадлежит более узкому классу функций. Например,


$\varphi(t_*)=\varphi_{0} \in R^n$, $\varphi(t) \equiv 0$, $t \in


[t_*-\beta, t_*]$. \tts{[} В таких случаях задача синтеза наблюдателя


упрощается.


\end{notes}





\begin{notes}


При длительных процессах наблюдения можно использовать


скользящее наблюдение [10], при котором процессы наблюдения


ведутся на промежутках $[\tau-\mu, \tau]$ заданной


продолжительности. Необходимые изменения в описанные выше


процедуры очевидны.


\end{notes}





\begin{notes}


Совокупность $x_{i}(t)$, $i=\overline{-1, M-1}$, не обязательно


строить в каждый момент $\tau \in T_{h}$. В предельном случае (см.


выше) можно найти ее только один раз.


\end{notes}





{\bf 2.2.} Рассмотрим теперь проблему наблюдения систем с


запаздыванием в условиях постоянно действующих возмущений


$$


\Dot x(t)=A(t)x(t)+A_{1}(t)x(t-\beta)+d(t)w(t).


$$


Здесь $w(t)$, $t \ge t_{*}$, --- неизвестное скалярное ограниченное


возмущение.





Задача синтеза наблюдателей для обыкновенных систем с возмущениями


исследована в [8]. Там она решена в классе возмущений,


генерируемых обыкновенными дифференциальными уравнениями с


неизвестным начальным состоянием. В этом случае задача сводилась к


стандартной задаче наблюдения состояния обыкновенной динамической


системы.





Если использовать приведенные выше результаты по наблюдению систем


с запаздыванием, то, следуя схеме [8], можно расширить класс


возможных возмущений и в качестве возмущений рассматривать решения


(выходные сигналы) нестационарных систем с запаздыванием с


неизвестной функцией $\varphi(\cdot)$ начального состояния


$$


w(t)=q'(t)y(t),\ \ \Dot y(t)=D(t)y(t)+D_{1}(t)y(t-\beta),


\ \ y_{t_*}(\cdot)=\psi_{t_*}(\cdot).


$$





\section{Стабилизация систем с запаздыванием по выходным сигналам}





{\bf 3.1.} Рассмотрим систему управления с запаздыванием в фазовых


переменных


\begin{equation}\label{31}


\Dot x(t)=A(t)x(t)+A_1(t)x(t-\beta)+B(t)u(t).


\end{equation}


Здесь $x(t)$, $u(t)$ --- $n$-векторы, $A(t)$, $A_1(t)$, $B(t)$, $t \ge


t_*$, --- ограниченные кусочно-непрерывные $n \times n$-матричные


функции, $\beta$, $0 < \beta < \infty$, --- запаздывание, $\det B(t)


\ge \alpha > 0$, $t \ge t_*$.





Пусть начальное состояние $x_{t_*}(\cdot)=\varphi_{t_*}(\cdot)$


неизвестно, но известно ограниченное множество $X_{t_*}(\cdot)$,


которому оно принадлежит. Информация о поведении $x(t)$, $t \ge


t_{*}$, системы (\ref{31}) в процессе управления поступает от


измерительного устройства (\ref{22}).





Предположим, что $y_t[\cdot]=y(s)$, $t _* \le s \le \tau$, ---


информация, записанная к моменту $\tau$, $Y_{t}[\cdot]$ --- множество


возможных сигналов $y_t[\cdot]$.





\begin{defns}


Функционал


\begin{equation}\label{32}


u=u(t,y_t[\cdot]),\ \ y_t \in Y_t[\cdot],


\ \ t \ge t_*,\ \ u(t,0)=0,\ \ t \ge t_*,


\end{equation}


называется стабилизирующей обратной связью по выходу, если нулевое


решение $x(t)=0$, $t \ge t_*$, замкнутой системы


\begin{equation}\label{33}


\Dot x(t)=A(t)x(t)+A_1(t)x(t-\beta)+B(t)u(t,y_t[\cdot]),


\ \ y(t)=c'(t)x(t),


\end{equation}


асимптотически устойчиво в $X_{t_*}(\cdot)$.


\end{defns}





Один из практических способов построения стабилизирующих обратных


связей по выходу (\ref{32}), (\ref{33}) состоит во включении в


канал $y_t[\cdot] \rightarrow u$ наблюдателя, задача которого ---


строить по выходным сигналам $y_t[\cdot]$ оценки $z_t(\cdot)$


состояний $x_t(\cdot)$, используемые далее регулятором для


выработки управляющего воздействия $u=u(t)=u(t,z_t(\cdot))$. Такой


подход отличается от метода создания гарантирующих управлений


[11], в котором оценки состояний не строятся, а доступная


информация $y_t[\cdot]$ используется непосредственно для


вычисления гарантированных оценок функционалов множеств,


необходимых для точного решения задачи оптимального управления.





Следуя [8], динамическую систему управления


\begin{equation}\label{34}


\Dot z (t) = A(t)z(t)+A_1(t)z(t-\beta)+D(t)v(t),


\ \ z_{t_*}(\cdot)=z_0(\cdot),


\end{equation}


назовем базой наблюдателя, если найдется такая обратная связь


$$


v=v(t,z_t(\cdot), y_t[\cdot]),\ \ z_t(\cdot) \in Z(\cdot),\ \ y_t[\cdot]


\in Y_t[\cdot],\ \ t \ge t_*,


$$


при которой решения $z(t)$, $t \ge t_*$, замкнутой системы


(наблюдателя)


$$


\Dot z (t) = A(t)z(t)+A_1(t)z(t-\beta)+D(t)v(t,z_t(\cdot), y_t[\cdot]),


\ \ z_{t_*}(\cdot)=z_0(\cdot),


$$


обладают свойством


\begin{equation}\label{35}


z_t(\cdot) \rightarrow x_t(\cdot),\ \ t \rightarrow \infty,


\end{equation}


при всех $y_t[\cdot] \in Y_{t}[\cdot]$, $t \ge t_*$.





В (\ref{34}) \ $D(t)$, $\det D(t) \ge \alpha > 0$, $t \ge t_*$, ---


ограниченная кусочно-непрерывная $n \times n$-матричная функция,


предел в (\ref{35}) понимается в смысле нормы $\|x_t(\cdot)\| =


\max \|x(s)\|$, $t -\beta \le s \le t$, $\|x\|=\max|x_i|$,


$i=\overline{1,n}$.





Будем говорить, что начальное состояниe $z_0(\cdot)$ и обратные связи


\begin{equation}{\label{36}}


u=u(t,z_t(\cdot)),\ \ v=v(t,z_t(\cdot),y_t[\cdot]),\ \ z_t(\cdot) \in


Z_t(\cdot),\ \ y_t[\cdot] \in Y_{t}[\cdot],\ \ t \ge t_*,


\end{equation}


реализуют стабилизирующую обратную связь по выходу (\ref{22}),


(\ref{32}), если решения $(x_t(\cdot),z_t(\cdot))$, $t \ge t_*$,


системы уравнений


\begin{gather}


\Dot x(t)=A(t)x(t)+A_1(t)x(t-\beta)+B(t)u(t,z_t(\cdot)),


\ \ x_{t_*}(\cdot) = \varphi_{t_*},\notag \\


%


\label{37}


\Dot z (t) =


A(t)z(t)+A_1(t)z(t-\beta)+D(t)v(t,z_t(\cdot),y_t[\cdot]),


\ \ z(t_0)=z_0;\ \ y(t)=c'(t)x(t),


\end{gather}


обладают свойством


\begin{equation}\label{38}


x(t) \rightarrow 0,\ \ t \rightarrow \infty,


\end{equation}


для всех $\varphi_{t_*}(\cdot) \in X_{t_*}$.





При конструировании систем управления наиболее часто встречающиеся


нелинейности связаны с насыщением. Поэтому большое значение с


точки зрения приложений имеют ограничения


$$


%\begin{equation}\label{39}


\|u(t,z_t(\cdot))\| \le L, \| v(t,z_t(\cdot), y_t[\cdot])\| \le


M, z_t(\cdot) \in Z_t(\cdot),\ \ y_t[\cdot] \in Y_t[\cdot],\ \ t \ge


t_*.


%\end{equation}


$$


Обратные связи, удовлетворяющие этим ограничениям, будем называть


ограниченными.





Как и выше, будем считать, что обратная связь может менять значение


выходных сигналов только в дискретные моменты времени $t \in


T_h=\{t_*, t_*+h, t_*+2h, \dotsc\}$. Тогда траектории нелинейных


уравнений (\ref{37}), замкнутых дискретными обратными связями


(\ref{36}), совпадают с траекториями линейных уравнений


(\ref{31}), (\ref{34}) с теми же начальными состояниями, но


находящимися под действием программных управлений (реализаций


обратных связей (\ref{36}))


\begin{gather*}


u(t)=u(t_*+kh, z_{t_*+kh}(\cdot)),


\ \ v(t)=v(t_*+kh,z_{t_*+kh}(\cdot),\ \ y_{t_*+kh}[\cdot]),\\


%


t \in [t_*+kh,t_*+(k+1)h], \tts{[} \ \ k=0,1,2,\dotsc


\end{gather*}





Переход к дискретным управлениям делает естественным замену


свойства (\ref{38}) на новое


$$


\|x(t)\| \rightarrow o(h),\ \ t \rightarrow \infty,


$$


которое при малых $h$ (что предполагается) с практической точки


зрения равносильно (\ref{38}).


\medskip





{\bf 3.2.} Пусть $\tau \in T_h$ --- текущий момент времени, $\tau


\ge (N+1)h$, $Nh=\beta$; $z_{\tau}^*(\cdot)$ --- текущее состояние


наблюдателя (\ref{34}), $y_{\tau}[\cdot]$ --- сигнал устройства


(\ref{22}), записанный к моменту $\tau$, $\mu_i(\tau)$, $i=\overline{0,Nn}$,


$0=\mu_0(\tau)<\mu_1(\tau)<\dotsb<\mu_{Nn}(\tau)$, --- некоторые


числа, кратные $h$; $0<\theta_u=N_u h<\infty$, $\theta_v=N_v h <


\infty$ --- параметры метода. Допустим, что выполняются условия


наблюдаемости


$$


\rank \left(


\begin{array}{c}


c'(\tau-\mu_{i}(\tau))F(\tau-\mu_i(\tau), t_*),


\ \ \displaystyle \int_{t_*-(k+1)h}^{t_*-kh}


c'(\tau-\mu_i(\tau)) \times \\[3mm] \times


F(\tau-\mu_i(\tau),s+\beta)


A_1(s+\beta)ds,\ \ k=\overline{0,N-1};


\ \ i=\overline{0,Nn}


\end{array} \right)=(N+1)n,


$$


где $F(t,s)$, $t \ge s$, --- фундаментальная $n \times n$-матрица


решений уравнения (\ref{31}) с $u(t)=0$, $t \ge t_*$; $x_{k}(\tau) \in


R^n$, $k=\overline{-1, N-1}$, --- решение системы


\begin{gather*}


c'(\tau-\mu_{i}(\tau))F(\tau-\mu_i(\tau), t_*)


+\sum^{N-1}_{k=0}


\int_{t_*-(k+1)h}^{t_*-kh}c'(\tau-\mu_i(\tau))\times \\


%


\times F(\tau-\mu_i(\tau),s+\beta)


A_1(s+\beta)ds\, x_k=y(\tau-\mu_i(\tau)),\ \ i=\overline{0,Nn}, \\


%


\wh{x}(t)=F(t,t_*)x_{-1}(\tau)+\sum_{k=0}^{N-1}


\int_{t_{*}-(k+1)h}^{t_{*}-kh}F(t,s+\beta)A_1(s+\beta)


ds\, x_k(\tau),\ \ t \ge t_*.


\end{gather*}





В классе дискретных управлений $v(t)$, $t \in [\tau,\tau+\theta_v]$, \tts{[}


кусочно-непрерывных управлений $v(t)$, $t \in [\tau+\theta_v,


\tau+\theta_v+\beta]$, рассмотрим задачу оптимального управления,


сопровождающую процесс наблюдения


\begin{gather}


B_v(\tau,z_\tau^*(\cdot),y_\tau[\cdot])=\min


\int_{\tau}^{\tau+\theta_v+\beta}\|v(t)\|dt, \notag \\


%


\label{310}


\Dot


z(t)=A(t)z(t)+A_1(t)z(t-\beta)+D(t)v(t),


\ \ z_{\tau(\cdot)}=z_{\tau}^*(\cdot),\\


%


D(t)v(t)+A_1(t)z(t-\beta)-A_1(t)\wh{x}(t-\beta)=0,\ \ t \in


[\tau+\theta_v,\tau+\theta_v+\beta], \notag \\


%


z(\tau+\theta_v)=\wh{x}(\tau+\theta_v),\ \ \|v(t)\| \le M,\ \ t \in


[\tau, \tau+\theta_v+\beta].\notag


\end{gather}





Обозначим через $v^0(t\mid\tau,z_{\tau}^*(\cdot)$, $y_\tau[\cdot])$, $t \in


[\tau,\tau+\theta_v+\beta]$, оптимальное программное решение


задачи (\ref{310}), $Z_\tau(\cdot)$ --- множество всех состояний


$Z_{\tau}^*(\cdot)$, для которых задача $(\ref{310})$ имеет решение


при всех $y_\tau[\cdot] \in Y_{\tau}[\cdot]$.





\begin{defns}


Функция


\begin{gather}\label{311}


v^0(\tau,z_{\tau},z_{\tau}^*(\cdot),y_\tau[\cdot])=


v^0(\tau\mid\tau,z_{\tau}^*(\cdot),y_\tau[\cdot]), \\


%


z_{\tau}^*(\cdot) \in Z_{\tau}(\cdot),\ \ y_{\tau}[\cdot] \in


Y_{\tau}[\cdot],\ \ \tau \in T_h,\notag


\end{gather}


называется оптимальным стартовым управлением типа дискретной


обратной связи задачи (\ref{310}).


\end{defns}





Алгоритмы реализации обратной связи (\ref{311}) аналогичны


описанным в [7].





Положим


\begin{equation}\label{312}


v(\tau,z_{\tau}(\cdot),y_{\tau}[\cdot])=


v^0(\tau,z_{\tau}(\cdot),y_{\tau}[\cdot]),


\ \ z_{\tau}(\cdot) \in Z_{\tau}(\cdot),\ \ y_{\tau}[\cdot] \in


Y_{\tau}[\cdot],\ \ \tau \in T_{h},


\end{equation}


возьмем $z_{0}(\cdot) \in Z_{t_*}(\cdot)$ и подадим сигнал


$v(t)=v(\tau,z_{\tau}^*(\cdot),y_{\tau}[\cdot])$, $t \in [\tau,


\tau+h]$, \tts{[}


на вход наблюдателя (\ref{34}). В результате в момент $\tau+h$


наблюдатель перейдет в состояние $z_{\tau+h}^*(\cdot)$.





В классе дискретных управлений $u(t)$, $t \in [\tau,


\tau+\theta_u]$, \tts{[} и кусочно-непрерывных управлений $u(t)$,


$t \in[\tau+\theta_u, \tau+\theta_u+\beta]$, рассмотрим задачу


оптимального управления, сопровождающую процесс стабилизации


\begin{gather}


B_{u}(\tau,z_{\tau}^*(\cdot))=


\min\int_{\tau}^{\tau+\theta_u+\beta}\|u(t)\|dt, \notag \\


%


\label{313}


\Dot x(t)=A(t)x(t)+A_1(t)x(t-\beta)+B(t)u(t),


\ \ x_{\tau}(\cdot)=z_{\tau}^*(\cdot), \\


%


B(t)u(t)+A_1(t)x(t-\beta)=0,\ \ t \in


[\tau+\theta_u,\tau+\theta_u+\beta], \notag \\


%


x(\tau+\theta)=0,


\ \ \|u(t)\|\le L,\ \ t \in [\tau,\tau+\theta_v+\beta].\notag


\end{gather}





Пусть $u^0(\tau,z_\tau^*(\cdot))$, $t \in


[\tau,\tau+\theta_u+\beta]$, --- оптимальное программное управление


задачи (\ref{313}), $X_{\tau}(\cdot)$ --- множество всех


$z_{\tau}^*(\cdot)$, для которых задача (\ref{313}) имеет решение.





Функцию


\begin{equation}\label{314}


u^0(\tau,z_\tau(\cdot))=u^0(\tau\mid\tau,z_{\tau}(\cdot)),


\ \ z_{\tau}(\cdot)\in X_{\tau}(\cdot),\ \ \tau \in T_h,


\end{equation}


назовем оптимальным стартовым управлением типа дискретной обратной


связи задачи (\ref{313}).





Реализация (\ref{314}) аналогична реализации (\ref{312}).





Положим


\begin{equation}\label{315}


u(\tau,z_{\tau}(\cdot))=u^0(\tau,z_{\tau}(\cdot)),\ \ z_{\tau}(\cdot)


\in X_{\tau}(\cdot),\ \ \tau \in T_h,


\end{equation}


и управление


$$


u^*(t)=u(\tau,z_{\tau}^*(\cdot)),\ \ t \in [\tau,\tau+h], \tts{[}


$$


подаем на вход системы (\ref{31}). После этого система (\ref{31})


перейдет в неизвестное положение $x^*(\tau+h)$, которое породит


доступный в момент $\tau+h$ сигнал $y(\tau+h)$. Этот сигнал вместе


с $y_\tau[\cdot]$ образует сигнал $y_{\tau+h}[\cdot]$.





Таким образом, для произвольного текущего момента времени описан


шаг процесса стабилизации системы с запаздыванием (\ref{31}) по выходу


(\ref{22}).





Можно при весьма общих условиях показать, что построенные обратные


связи (\ref{312}), (\ref{315}) и начальное состояние


$z_{0}(\cdot)$ реализуют ограниченную стабилизирующую дискретную


обратную связь по выходу.


Стабилизация динамических


систем с запаздыванием в управлении (см. п.\,1.5) по выходным


сигналам (\ref{22}) осуществляется с помощью наблюдателя обыкновенных


систем [8] и стабилизатора по состоянию для систем с запаздыванием


в управлении (см. п.\,1.1).








\begin{thebibliography}{99}





\bibitem{1}


Габасов Р., Кириллова Ф.М. и др. {\it Конструктивные методы


оптимизации}. Ч.\,5. {\it Нелинейные задачи}. -- Минск:


Изд-во ``Университетское'', 1998. -- 501 с.


\bibitem{2}


Красовский Н.Н. {\it О стабилизации динамических систем


дополнительными силами} // Дифференц. уравнения. -- 1965. --


Т.\,1. -- \No\,1. -- С.\,5--16.


\bibitem{3}


Осипов Ю.С. {\it О стабилизации управляемых систем с


запаздыванием} // Дифференц. уравнения. -- 1965. -- Т.\,1. --


\No\,5. -- С.\,605--618.


\bibitem{4}


Габасов Р., Кириллова Ф.М. {\it Стабилизация динамических


систем методами оптимального управления} // Докл. НАН Беларуси. --


1998. -- Т.\,42. --  \No\,1. -- С.\,18--23.


\bibitem{5}


Габасов Р., Ружицкая Е.А. {\it Стабилизация динамических систем


с обеспечением дополнительных свойств переходных процессов} //


Кибернет. и системн. анализ. -- 2001. -- \No\,3. -- C.\,139--151.


\bibitem{6}


Габасов Р., Кириллова Ф.М. {\it Стабилизация нелинейных систем


при больших начальных возмущениях}


// Докл. НАН Беларуси. -- 2001. -- Т.\,45. -- \No\,4. -- С.\,5--8.


\bibitem{7}


Балашевич Н.В., Габасов Р., Кириллова Ф.М. {\it Численные


методы программной и позиционной оптимизации линейных систем


управления} // Журн. вычисл. матем. и матем. физ. -- 2000. -- 


Т.\,40. -- \No\,6. -- С.\,838--859.


\bibitem{8}


Габасов Р., Кириллова Ф.М. {\it Синтез нелинейных наблюдателей


динамических систем методами оптимального управления}


// Докл. НАН Беларуси. -- 2002. -- Т.\,46. -- \No\,4. -- С.\,5--8.


\bibitem{9}


Балашевич Н.В., Габасов Р., Кириллова Ф.М. {\it Алгоритмы


программной и позиционной оптимизации систем управления с


промежуточными фазовыми ограничениями}


// Журн. вычисл. матем. и матем. физ. -- 2001. -- Т.\,41.


-- \No\,10. -- С.\,1485--1504.


\bibitem{10}


Габасов Р., Дмитрук Н.М., Кириллова Ф.М. {\it Скользящее


наблюдение за поведением динамических систем}


// Докл. НАН Беларуси. -- 2001. -- Т.\,45. -- \No\,2. -- С.\,12--16.


\bibitem{11}


Габасов Р., Кириллова Ф.М. {\it Конструктивный метод


гарантированной оптимизации нелинейных динамических систем} //


Докл. НАН Беларуси. -- 2000. -- Т.\,44. -- \No\,4. -- С.\,5--8.





\end{thebibliography}





\vskip 2cm





\post{\it Белорусский государственный университет}{\it Поступила}


\post{\it Институт математики Национальной}{02.09.2002}


\post{\it Академии наук Беларуси}{}





\label{end}


\end{document}





